\section{Exceptional labels and opposite coalescence}
\label{sec:exceptional}

Fix one of the charts of Lemma~\ref{lem:charts}, with upper label
set $I$ and opposite label set $-I$.  Throughout the deterministic
arguments below, work in an environment satisfying the simultaneous
conclusions of Theorem~\ref{thm:structure}, Lemma~\ref{lem:uniform-metric},
and Proposition~\ref{prop:cone}.  The coalescence assertion in
Theorem~\ref{thm:structure} is only a statement for each fixed label.
We shall obtain the following different statement: every label at which
branching or Busemann nonuniqueness occurs has a coalescing family at its
opposite label.  The exceptional label may depend on the environment.

\Needspace{6\baselineskip}
\subsection{Calibrated rays and countability}

\begin{lemma}[Rays calibrated from an arbitrary vertex]
\label{lem:calibrated-rays}
Let $p$ be a ray, let $B=B_p$, and let $x\in\Z^2$.  There is a ray
$q=(q_j)_{j\ge0}$ from $x$ such that
\begin{equation}
  T(x,q_j)=B(x,q_j)\qquad(j\ge0).
  \label{eq:calibrated-from-x}
\end{equation}
Moreover, $q$ has the same label as $p$.  In general the conclusion
concerns the labels; it does not assert $B_q=B_p$.
\end{lemma}

\begin{proof}
Take the finite geodesics from $x$ to $p_k$.  Since $p$ is simple,
$|p_k|\to\infty$, and their lengths tend to infinity.  Local finiteness
and a diagonal subsequence give a local limit $q$ from $x$.  Each finite
prefix of $q$ is eventually a prefix of these finite geodesics, so $q$
is a simple infinite geodesic.

By Lemma~\ref{lem:busemann-basic},
\begin{equation}
 d_k(x):=T(x,p_k)-B(x,p_k)\ge0,
 \qquad d_k(x)\longrightarrow0.
 \label{eq:calibrated-limit-defect}
\end{equation}
If $w$ lies on the geodesic from $x$ to $p_k$, then
\[
 0\le T(x,w)-B(x,w)
 \le T(x,p_k)-B(x,p_k)=d_k(x),
\]
because $T(w,p_k)\ge B(w,p_k)$.  For each fixed vertex of the local
limit, let $k$ tend to infinity along the chosen subsequence.  This
proves \eqref{eq:calibrated-from-x}.

For any $w\in\Z^2$, the same metric bound and
\eqref{eq:calibrated-from-x} give
\[
 T(x,q_j)-T(w,q_j)
 \le B(x,q_j)-B(w,q_j)=B(x,w).
\]
Consequently $B_q(x,w)\le B(x,w)$.  Write $\sigma=\rho_q$ and
$\rho=\rho_p$.  Asymptotic linearity, with the fixed basepoint $x$,
and the choices $w=ke_i$ and $w=-ke_i$, $i=1,2$, imply respectively
$\sigma(e_i)\le\rho(e_i)$ and $\sigma(e_i)\ge\rho(e_i)$.
Thus $\sigma=\rho$.
\end{proof}

Call a label $\rho\in I$ \emph{exceptional} in an environment $\omega$
if at least one of the following holds:
\begin{enumerate}[label=(\roman*)]
\item two distinct rays of label $\rho$ start at the same vertex;
\item two rays of label $\rho$ have different Busemann functions.
\end{enumerate}
Denote this set by $\mathcal E_I(\omega)$.

\begin{proposition}[Countability of exceptional labels]
\label{prop:exceptional-countable}
Almost surely $\mathcal E_I(\omega)$ is at most countable.  The same
conclusion holds with $I$ replaced by $-I$.
\end{proposition}

\begin{proof}
We prove the two assertions in the definition separately.

\emph{Branching.}
Fix a vertex $x$ and two different first edges $e,f$ incident to $x$.
There are at most two labels in $I$ that are realized both by a ray
from $x$ beginning with $e$ and by one beginning with $f$.  Suppose,
to the contrary, that there are three such labels
$\rho_1,\rho_2,\rho_3$, with
$r(\rho_1)<r(\rho_2)<r(\rho_3)$.  Choose the six corresponding rays.
The geodesics from $x$ form a tree: once two of them diverge, they
cannot meet again, by uniqueness of finite geodesics.  In particular
the two first-edge branches are disjoint away from $x$.

Choose a horizontal row sufficiently high that its six first-arrival
vertices are distinct and have the label order supplied by
Lemma~\ref{lem:ray-order}.  Such a row exists: rays of different
labels have disjoint tails, and rays in the two first-edge branches
never meet away from $x$.  The endpoints occur as three consecutive
groups, one for each label, and each group contains one endpoint of
each branch.  Color the endpoints according to the first edge $e$ or
$f$.  The resulting six-letter word is a concatenation of three
two-letter words, each containing both colors.  It has at least three
changes of color, and hence contains an alternating four-letter
subsequence.

Connect the two $e$-colored endpoints of this subsequence by the
unique arc in the $e$ branch, and connect the two $f$-colored endpoints
by the corresponding arc in the $f$ branch.  Both arcs avoid $x$,
since the two root paths defining each arc share their first edge.
Their interiors are strictly below the chosen row, since all six
root paths were stopped on their first visit to that row.  The arcs
are disjoint, but their endpoints alternate along the row.  This is
impossible: the first arc together with the row segment between its
ends is a Jordan curve, and the second arc must cross it away from
the row.  This proves the claim for $x,e,f$.

Two distinct rays from a common vertex have a first vertex of
divergence.  Applying the preceding claim there, and taking the
countable union over vertices and first-edge pairs, proves
countability for condition~(i).

\emph{Busemann nonuniqueness.}
First consider two vertices $x,y$, with $y$ sufficiently far to the
right of $x$ that the crossing conclusion of
Lemma~\ref{lem:crossing} applies to every pair of upper rays from
these vertices with reversed strict label order.  If $B,B'$ are
Busemann functions of upper rays with labels $\rho,\sigma$ satisfying
$r(\rho)<r(\sigma)$, then
\begin{equation}
 B(x,y)\le B'(x,y).
 \label{eq:ordered-busemann-values}
\end{equation}
Indeed, use Lemma~\ref{lem:calibrated-rays} to take a $B$-calibrated
ray from $y$ and a $B'$-calibrated ray from $x$.  Their labels are
$\rho$ and $\sigma$, respectively, so they intersect at a vertex
$w$.  Calibration at $w$ and the metric bounds on Busemann increments
give
\begin{align*}
 B(x,y)&=B(x,w)-T(y,w)\le T(x,w)-T(y,w)\\
       &\le T(x,w)-B'(y,w)=B'(x,y).
\end{align*}

For such a pair $x,y$, let
\[
 S_\rho(x,y)=\{B_p(x,y):p\text{ is a ray with label }\rho\}.
\]
These sets are bounded by $T(x,y)$ in absolute value.  If
$r(\rho)<r(\sigma)$ and both sets are nonempty,
\eqref{eq:ordered-busemann-values} states that
\begin{equation}
 \sup S_\rho(x,y)\le\inf S_\sigma(x,y).
 \label{eq:ordered-value-ranges}
\end{equation}
Thus the nonempty open intervals
$(\inf S_\rho(x,y),\sup S_\rho(x,y))$ are pairwise disjoint.  Each
contains a rational number, so only countably many labels can have
more than one possible value at this pair.

It remains to ensure that these comparisons detect every possible
disagreement of Busemann functions.  Fix arbitrary vertices $u,v$.
Proposition~\ref{prop:cone} supplies finite offsets controlling all
upper rays from $u$ and from $v$.  For $y_N=(N,0)$ it also supplies
offsets $H_N=o(N)$ controlling all upper rays from $y_N$.  The height
differences between $y_N$ and $u,v$ are fixed.  For all sufficiently
large integers $N$, the horizontal separations therefore dominate
the offsets and height differences as required by
Lemma~\ref{lem:crossing}.  Choose any such $N$.  Both pairs $(u,y_N)$
and $(v,y_N)$ then satisfy \eqref{eq:ordered-value-ranges} for all
labels and all their Busemann functions.

If two upper-ray Busemann functions $B,B'$ of the same label disagree
at $(u,v)$, additivity in
\[
 B(u,v)=B(u,y_N)-B(v,y_N),\qquad
 B'(u,v)=B'(u,y_N)-B'(v,y_N)
\]
forces a disagreement at one of those two comparison pairs.  The
labels producing a disagreement at $(u,v)$ are consequently contained
in a union of two countable sets.  Taking a countable union over
$u,v\in\Z^2$ proves countability for condition~(ii).  The choice of
$N$ may depend on the environment and on $u,v$; it uses the all-ray
cone bounds, and no common rate of Busemann convergence over labels.
Reflection of the vertical coordinate proves the lower-sector
assertion.
\end{proof}

\subsection{Independent environments and exceptional parameters}

The probabilistic use of countability requires care because
$\mathcal E_I$ is random.  We give the measurable argument before
constructing the coupling that will use it.

\begin{lemma}[Countable exceptional parameters in an independent metric]
\label{lem:independent-exceptional}
Let $V,W$ be independent environments, each with the iid law $\PP$.
Almost surely, simultaneously for every $\rho\in\mathcal E_I(V)$,
all $W$-rays of label $-\rho$ coalesce.
\end{lemma}

\begin{proof}
Use the standard Borel configuration space $\Omega$ before completing
its probability measure, and the path space
$\mathcal P=(\Z^2)^{\mathbb N_0}$.  Let
$\mathsf R(\omega,p,\rho)$ be the assertion that $p$ is a simple
nearest-neighbor ray in $\omega$, that all its Busemann limits exist,
and that its Busemann function is asymptotically linear to $\rho$.
For a linear functional $\rho$, this is a Borel relation by
Lemma~\ref{lem:measurable}. Define Busemann evaluations to be zero
where their limits do not exist; the relation $\mathsf R$ explicitly
requires existence.

For paths $p_1,p_2$ satisfying $\mathsf R(\omega,p_i,\rho)$, the
exceptionality witness is the Borel disjunction
\begin{equation}
 \bigl[p_{1,0}=p_{2,0}\ \text{and}\ p_1\ne p_2\bigr]
 \quad\text{or}\quad
 \bigl[\exists x,y\in\Z^2:
 B_{p_1}(x,y)\ne B_{p_2}(x,y)\bigr].
 \label{eq:exceptional-witness-test}
\end{equation}
The common-tail test for paths $q_1,q_2$ is Borel as well:
\begin{equation}
 \exists a,b\in\mathbb N_0\quad
 \forall k\in\mathbb N_0:\ q_{1,a+k}=q_{2,b+k}.
 \label{eq:common-tail-test}
\end{equation}
Its negation tests noncoalescence.

Now define the bad set $J\subset\Omega\times\Omega$ by the existence
of one label $\rho\in I$ and four paths $p_1,p_2,q_1,q_2$ such that
\begin{enumerate}[label=(\alph*)]
\item $\mathsf R(v,p_i,\rho)$ holds for $i=1,2$, and
      \eqref{eq:exceptional-witness-test} holds in $v$;
\item $\mathsf R(w,q_i,-\rho)$ holds for $i=1,2$, and
      \eqref{eq:common-tail-test} fails.
\end{enumerate}
This is directly the projection of a Borel relation on
$\Omega^2\times I\times\mathcal P^4$.  In particular, $J$ is
analytic and is measurable in the completion of $\PP\otimes\PP$.

For every $v$ in the full-measure countability event of
Proposition~\ref{prop:exceptional-countable}, its section is
\begin{equation}
 J_v=\bigcup_{\rho\in\mathcal E_I(v)}
 \{w:\text{two $w$-rays of label $-\rho$ do not coalesce}\}.
 \label{eq:countable-null-section}
\end{equation}
For this fixed $v$, each label in the countable union is an ordinary
fixed parameter.  Symmetry gives $-\rho\in C$, and
Theorem~\ref{thm:structure} makes the corresponding analytic event
in $w$ null.  Hence $\PP(J_v)=0$ for almost every $v$.
Completed-product Fubini gives
$(\PP\otimes\PP)(J)=0$.  Equivalently, replace $J$ by a Borel set
equal to it outside a Borel product-null set and apply ordinary
Fubini; the sections of that null set are null for almost every
$v$.  This argument needs no measurable enumeration of
$\mathcal E_I(v)$.

No universal all-ray condition is inserted into the four-path
relation. Such regularity conditions are imposed outside the
projection, as specified after Lemma~\ref{lem:measurable}.
\end{proof}

\subsection{Transfer across a horizontal cut}

Let $U$ have the iid law.  For $n\ge1$, set
\[
 H_n=\{z\in\Z^2:z_2=-n\},\qquad
 E_n^+=\{\{x,y\}:x\sim y,\ x_2\ge-n,\ y_2\ge-n\}.
\]
Take independent iid replacement arrays $A^{(n)},D^{(n)}$, all
independent of $U$, and define
\begin{equation}
 (V_n)_e=
 \begin{cases}U_e,&e\in E_n^+,\\ A^{(n)}_e,&e\notin E_n^+,\end{cases}
 \qquad
 (W_n)_e=
 \begin{cases}D^{(n)}_e,&e\in E_n^+,\\ U_e,&e\notin E_n^+.\end{cases}
 \label{eq:half-plane-coupling}
\end{equation}
For each $n$, the pair $(V_n,W_n)$ has law $\PP\otimes\PP$: the
coordinates of $U$ retained in its two members are disjoint, and
their replacement coordinates are independent.  Independence of
the pairs for different $n$ is neither asserted nor needed.

We will use $V_n$ to retain the two upper rays witnessing an exceptional
label, with their exact label and any distinguishing Busemann evaluation.
After choosing one such cut, $W_n$ must retain tails of arbitrary rays
with the opposite label, so that Lemma~\ref{lem:independent-exceptional}
can be applied to the independent pair $(V_n,W_n)$.

\begin{lemma}[Half-plane transfer]
\label{lem:cut-transfer}
There is a deterministic increasing sequence $n_k\to\infty$ and an
event of probability one for the coupling
\eqref{eq:half-plane-coupling} with the following properties.
\begin{enumerate}[label=(\roman*)]
\item For every finite collection of upper $U$-rays and every finite
      set $K\subset\Z^2$, all sufficiently large $k$ have the
      following property.  Every ray $p$ in that collection is a
      $V_{n_k}$-ray with exactly its original label, and
      \[
       B_p^{V_{n_k}}(x,y)=B_p^U(x,y)\qquad(x,y\in K).
      \]
      More precisely, for each such $k$ and $p$, the distances
      $T_{V_{n_k}}(x,p_j)$ and $T_U(x,p_j)$ are equal for all
      sufficiently large $j$, simultaneously for $x\in K$.
\item For each fixed $n\ge1$, every lower $U$-ray has a tail which
      is a $W_n$-ray with exactly its original label.
\end{enumerate}
The thresholds may depend on the rays, their labels, $K$, and, in
part~(ii), on $n$.
\end{lemma}

\begin{proof}
First take a common event on which all-ray asymptotics, finite
geodesic existence and uniqueness, and the uniform metric estimate
hold in $U$ and in all $V_n,W_n$.  This is a countable intersection
of events of probability one.

\emph{A single subsequence for the moving cuts.}
For any metric $T$ put
\begin{equation}
 E_n(T)=\sup_{x,y\in\Z^2}
 \frac{|T(x,y)-m(y-x)|}{|x|+|y|+n}.
 \label{eq:moving-metric-error}
\end{equation}
For a fixed typical metric, $E_n(T)\to0$.  Indeed, for any
$\eps>0$, Lemma~\ref{lem:uniform-metric} bounds the numerator by
$\eps(|x|+|y|)$ outside a finite set of pairs.  The maximum error
on the remaining finite set, divided by $n$, tends to zero.  Since
each $V_n$ has the iid law, $E_n(T_{V_n})\to0$ in probability.
Choose an increasing deterministic sequence $n_k$ such that
\[
 \PP\bigl(E_{n_k}(T_{V_{n_k}})>2^{-k}\bigr)<2^{-k}.
\]
After a further probability-one restriction, Borel--Cantelli gives
\begin{equation}
 E_{n_k}(T_{V_{n_k}})\longrightarrow0.
 \label{eq:subsequence-metric-error}
\end{equation}
This is the only subsequence used in the proof.

Fix now an arbitrary upper $U$-ray $p$, with label $\rho\in I$.
Write
\begin{equation}
 a(w)=B_p^U(0,w)-\rho(w)=o(|w|),\qquad
 B_p^U(x,z)=\rho(z-x)+a(z)-a(x).
 \label{eq:rooted-busemann-error}
\end{equation}
All contacts of $\rho$ are strictly upper, so compactness gives
\begin{equation}
 \gamma=\min_{|d|=1,\ d_2\le0}\bigl(m(d)-\rho(d)\bigr)>0.
 \label{eq:opposite-semicircle-gap}
\end{equation}
Let $K$ be finite and enlarge it to contain $p_0$.  Set
$R_K=\max_{x\in K}|x|$.  For $n\ge2R_K+1$, $x\in K$, and
$z\in H_n$, the vector $z-x$ points into the lower half-plane and
\begin{equation}
 |z-x|\ge |z|/2,\qquad
 |x|+|z|+n\le \tfrac52|z|\le5|z-x|.
 \label{eq:moving-cut-comparability}
\end{equation}
For $n=n_k$, \eqref{eq:moving-metric-error} therefore yields
\begin{align}
 T_{V_n}(x,z)-B_p^U(x,z)
 &\ge \gamma|z-x|
      -E_n(T_{V_n})(|x|+|z|+n)-|a(z)|-|a(x)|.
 \label{eq:moving-cut-barrier-estimate}
\end{align}
The metric error is $o(|z-x|)$ uniformly over these $x,z$, by
\eqref{eq:subsequence-metric-error} and
\eqref{eq:moving-cut-comparability}.  The same is true of $|a(z)|$:
$|z|\ge n\to\infty$ uniformly on the cut, and
$|z|\le2|z-x|$.  Finally, the finitely many $|a(x)|$ are bounded.
Thus, for all sufficiently large $k$, uniformly for $x\in K$ and
$z\in H_{n_k}$,
\begin{equation}
 T_{V_{n_k}}(x,z)-B_p^U(x,z)
 \ge \frac\gamma2|z-x|\ge\frac\gamma4 n_k.
 \label{eq:moving-cut-positive-barrier}
\end{equation}
The same estimate holds with $V_{n_k}$ replaced by $U$, because
$E_n(T_U)\to0$.

\emph{The last-cut comparison.}
The heights of $p_j$ tend to $+\infty$, so the entire ray lies
strictly above $H_n$ for all sufficiently large $n$.  If a finite
walk from $x$ to $p_j$, with both endpoints above the cut, visits
$H_n$, let $z$ be its last cut vertex.  Its suffix after $z$ lies
strictly above the cut and uses only edges in $E_n^+$.  Its cost in
$V_n$ is consequently at least
\begin{align}
 T_{V_n}(x,z)+B_p^U(z,p_j)
 &=B_p^U(x,p_j)
   +\bigl[T_{V_n}(x,z)-B_p^U(x,z)\bigr].
 \label{eq:last-cut-lower-bound}
\end{align}
The suffix bound uses only $B_p^U(z,p_j)\le T_U(z,p_j)$ and its
unchanged edge costs.  Repeated visits to the cut and repeated edges
in the walk cause no problem: all occur either in the prefix or in
a suffix whose cost still bounds the corresponding distance.  The
same comparison holds in $U$.

For $x=p_0$, the original segment of $p$ has cost
$B_p^U(p_0,p_j)$ and is unchanged.  By
\eqref{eq:moving-cut-positive-barrier} every cut-touching competitor
is more expensive; a competitor avoiding the cut has its original
$U$ cost.  Hence every prefix of $p$ minimizes in $V_{n_k}$ for all
sufficiently large $k$, and $p$ is a ray there.

For any fixed $x\in K$, \eqref{eq:calibrated-limit-defect}, applied in
$U$, says that
\[
 T_U(x,p_j)-B_p^U(x,p_j)\longrightarrow0.
\]
Fix a sufficiently large $k$.  Eventually this defect is smaller
than $\gamma n_k/4$.  By the $U$ version of
\eqref{eq:last-cut-lower-bound}, a $U$-geodesic to $p_j$ must then
avoid the cut.  It provides an unchanged competitor in $V_{n_k}$.
The $V_{n_k}$ version excludes every cut-touching walk that could
improve that competitor, and every walk avoiding the cut has its
$U$ cost.  It follows that
\begin{equation}
 T_{V_{n_k}}(x,p_j)=T_U(x,p_j)
 \quad\text{for all sufficiently large }j.
 \label{eq:finite-witness-distance-equality}
\end{equation}
Because $K$ is finite, one target-index threshold works for all its
vertices.  Subtracting two such distance equalities and taking the
limit preserves every $B_p(x,y)$ with $x,y\in K$.

\emph{Exact identification of the label at a fixed cut.}
Preservation of a path alone does not identify both coordinates of
its Busemann asymptote.  Fix one $n=n_k$ for which $p$ was preserved,
and henceforth regard $U,V_n$, and the cut as fixed.  Fix $0<c<1$,
take $x$ in the sector $x_2\ge c|x|$, and let $z\in H_n$.  With
$d=z-x$ we have
\[
 |d|\ge x_2+n\ge c|x|,
 \qquad |z|\le|d|+|x|.
\]
Consequently
\begin{equation}
 |z-x|\ge\kappa(|x|+|z|),\qquad
 \kappa=\frac{c}{c+2}>0.
 \label{eq:fixed-cut-comparability}
\end{equation}
This bound holds uniformly along the entire infinite cut, including
cut vertices with arbitrarily large horizontal coordinate.  Since
$d_2<0$, \eqref{eq:opposite-semicircle-gap} gives
\[
 m(d)-\rho(d)\ge\gamma\kappa(|x|+|z|).
\]
For every $\eps>0$ there is $C_\eps<\infty$ such that
$|a(w)|\le\eps|w|+C_\eps$ for every $w$.  The fixed-metric uniform
estimate, separately in $U$ and $V_n$, and
\eqref{eq:rooted-busemann-error} now give, for $|x|$ sufficiently
large in the sector and every $z\in H_n$,
\begin{align*}
 T_{V_n}(x,z)-B_p^U(x,z)
 &\ge (\gamma\kappa-2\eps)(|x|+|z|)-2C_\eps.
\end{align*}
Choose $\eps<\gamma\kappa/8$ and then increase the lower bound on
$|x|$.  In both metrics we obtain the uniform barrier
\begin{equation}
 \min\{T_{V_n}(x,z),T_U(x,z)\}-B_p^U(x,z)
 \ge\frac{\gamma\kappa}{2}(|x|+|z|).
 \label{eq:fixed-cut-positive-barrier}
\end{equation}
In particular, the error at the moving basepoint $x$, as well as
that at $z$, has been included explicitly.

For each such $x$, its calibration defect tends to zero as the
target runs out along $p$.  The last-cut argument just used proves
$T_{V_n}(x,p_j)=T_U(x,p_j)$ for all sufficiently large $j$.
Meanwhile the preserved root $p_0$ satisfies
$T_{V_n}(p_0,p_j)=T_U(p_0,p_j)$ for every $j$, whether or not $p_0$
is in the strict sector.  Thus
\begin{equation}
 B_p^{V_n}(p_0,x)=B_p^U(p_0,x)
 \quad\text{for all sufficiently far }x\text{ with }x_2\ge c|x|.
 \label{eq:sector-busemann-equality}
\end{equation}
Take $c=1/2$, and use first $x=ke_2$ and then
$x=k(e_1+e_2)$.  Dividing \eqref{eq:sector-busemann-equality} by $k$
and using asymptotic linearity in the two fixed environments shows
that their labels agree on $e_2$ and $e_1+e_2$.  These vectors are
independent, so the labels are exactly equal.

All steps so far are pathwise on the common event already chosen.
For a finite collection of rays, apply them to each ray and take
the largest of the finitely many thresholds.  This proves
part~(i), with the stated quantifiers.

\emph{Lower tails for a fixed $W_n$.}
Fix $n$ and an arbitrary lower $U$-ray $q$ of label $\lambda\in-I$.
Use
\begin{gather*}
 a_-(w)=B_q^U(0,w)-\lambda(w)=o(|w|),\\
 \gamma_-=\min_{|d|=1,\ d_2\ge0}(m(d)-\lambda(d))>0.
\end{gather*}
Choose $c\in(0,1/2]$ small enough that eventually
$q_{j,2}\le-c|q_j|$.  This is possible because the limit directions
of $q$ lie in the strict lower sector of the chart.  If
$x_2\le-c|x|$, $|x|\ge2n/c$, and $z\in H_n$, then
\[
 |z-x|\ge c|x|-n\ge\tfrac c2|x|,
 \qquad |x|+|z|\le(1+4/c)|z-x|.
\]
The displacement $z-x$ points upward.  Applying the fixed-metric
uniform estimate to $U,W_n$ and bounding both errors
$a_-(x),a_-(z)$ exactly as above gives, for all sufficiently large
$|x|$ in this lower sector, uniformly for $z\in H_n$,
\begin{equation}
 \min\{T_{W_n}(x,z),T_U(x,z)\}-B_q^U(x,z)
 \ge\frac{\gamma_-c}{2(c+4)}(|x|+|z|).
 \label{eq:lower-fixed-cut-barrier}
\end{equation}
Only the usual uniform estimate in this particular fixed $W_n$ has
been used.

Choose a vertex $a=q_L$ sufficiently far along $q$ that
\eqref{eq:lower-fixed-cut-barrier} applies to $a$ and the whole
remaining tail lies strictly below $H_n$.  For a walk from $a$ to
a later $q_j$ that touches the cut, take its last cut vertex $z$.
After this last visit its first edge goes downward and every later
vertex is strictly below the cut.  Every edge of this suffix is
outside $E_n^+$ and hence is shared by $U$ and $W_n$.  Horizontal
edges on the cut, which are not shared, are part of the prefix.
The analogue of \eqref{eq:last-cut-lower-bound}, now using
\eqref{eq:lower-fixed-cut-barrier}, shows that this walk is more
expensive than the unchanged segment of $q$.  Walks that avoid the
cut have their $U$ costs.  Thus the tail from $a$ is a $W_n$-ray,
and its distances from $a$ to all later $q_j$ are unchanged.

Finally take any sufficiently far $x$ in the strict lower sector.
Its $U$ calibration defect relative to $q_j$ tends to zero, so
\eqref{eq:lower-fixed-cut-barrier} and the last-cut argument give
eventual equality of distances from $x$ in the two environments.
Using the preserved anchor $a$ yields equality of the two tail
Busemann functions at $(a,x)$.  The sequences $x=-ke_2$ and
$x=-k(e_1+e_2)$ belong to this sector for our choice $c\le1/2$.
Asymptotic linearity along them identifies the label of the
$W_n$-tail as exactly $\lambda$.

Since all $W_n$ were included in the initial countable
probability-one event, the argument applies to every fixed $n$ and
every lower ray, with its own tail index.  There is no estimate on
the rate of convergence of the sequence of metrics $W_n$, and no
uniformity of thresholds over labels is required.
\end{proof}

\subsection{The simultaneous conclusion}

\begin{proposition}[Coalescence at opposite exceptional labels]
\label{prop:opposite-coalescence}
Almost surely, simultaneously for all $\rho\in\mathcal E_I(U)$,
every two $U$-rays of label $-\rho$ coalesce.  Almost surely the same
statement holds with upper and lower sectors exchanged.  These
conclusions hold simultaneously for the countable collection of
charts under consideration.
\end{proposition}

\begin{proof}
Use the coupling \eqref{eq:half-plane-coupling}.  For every $n$,
Lemma~\ref{lem:independent-exceptional} applies to $(V_n,W_n)$.
Take their countable intersection, together with the
probability-one event in Lemma~\ref{lem:cut-transfer}.

On this event, fix any $\rho\in\mathcal E_I(U)$ and two upper rays
witnessing its exceptionality.  If they witness branching, their
distinct paths and common starting vertex are preserved by
Lemma~\ref{lem:cut-transfer}(i) for every sufficiently deep
subsequence cut.  If instead they witness different Busemann
functions, choose one pair $x,y$ at which those functions differ
and include $x,y$ in the finite witness set of that lemma.  At
every sufficiently deep subsequence cut the two paths retain both
their exact label $\rho$ and those unequal evaluations.  In either
case choose one such cut $n$; then
$\rho\in\mathcal E_I(V_n)$.

Now let $q,q'$ be any two lower $U$-rays of label $-\rho$.
Lemma~\ref{lem:cut-transfer}(ii), for this fixed $n$, supplies tails
of both which are $W_n$-rays of the same exact label $-\rho$.
They coalesce by the property of $(V_n,W_n)$.  Equality of their
tails is equality of their vertex sequences and thus also means
that the original $U$-rays coalesce.

The argument made arbitrary pathwise choices of $\rho$ and its
witnesses after the common event was fixed.  It therefore proves
the simultaneous assertion in the coupling, and hence for its
$U$ marginal.  Its failure is itself the analytic four-witness
event with both environments set equal to $U$, so it is measurable
in the completed marginal space.  Reflection gives the reversed
statement; a countable intersection gives both statements in all
charts.
\end{proof}
